% Vietnamese translation for OI Public Library
% Source-File: IOI中国国家候选队论文/2007/day1/8.杨哲《凸完全单调性的一个加强与应用》.ppt
\documentclass[11pt]{article}
\usepackage{oipl-vn}

\newcommand{\Oh}{\mathcal{O}}
\newcommand{\opt}{\operatorname{opt}}

\title{Một tăng cường của tính đơn điệu lồi hoàn toàn và ứng dụng}
\author{Yang Zhe}
\date{2007}

\begin{document}
\maketitle

\origtitle{A strengthening of convex total monotonicity and applications}
\sourcepath{IOI China candidate team papers / 2007 / day 1 / Yang Zhe.ppt}

\section*{Tóm tắt}

Bản trình chiếu bàn về một dạng mạnh hơn của tính đơn điệu lồi hoàn toàn
trong tối ưu quy hoạch động. Văn bản trích xuất được từ PowerPoint chứa các
công thức với hàm chi phí \(w(i,j)\), bất đẳng thức kiểu Monge, các hàm
\[
  t(i,x)=f[i]+w(i,x),
\]
biên \(B[i]\), và các ví dụ như CEOI 2004 \textit{Two Sawmills}. Bản ghi chú
này diễn giải mạch chính: từ điều kiện lồi rời rạc, suy ra thứ tự của các
điểm giao giữa ứng viên, rồi dùng thứ tự ấy để tính một lớp DP trong thời
gian tuyến tính hoặc gần tuyến tính.

\paragraph{Từ khóa.}
Tính đơn điệu hoàn toàn; bất đẳng thức Monge; tối ưu quy hoạch động; bao lồi
rời rạc; Two Sawmills.

\section{Dạng quy hoạch động}

Nhiều bài toán có chuyển trạng thái
\[
  f[x]=\min_{i<x}\{f[i]+w(i,x)\}.
\]
Nếu xét mọi \(i\) cho từng \(x\), thời gian là \(\Oh(n^2)\). Ta gọi
\[
  t(i,x)=f[i]+w(i,x)
\]
là giá trị do ứng viên \(i\) sinh ra tại vị trí \(x\). Bài toán trở thành:
với mỗi \(x\), tìm ứng viên \(i\) cho giá trị \(t(i,x)\) nhỏ nhất.

Ý tưởng của slide là nếu \(w\) có tính lồi rời rạc đủ mạnh, thì các ứng viên
tốt nhất xuất hiện theo thứ tự. Khi đó ta chỉ cần duy trì danh sách các ứng
viên còn có thể tối ưu và biên bắt đầu có hiệu lực của từng ứng viên.

\section{Điều kiện lồi và Monge}

Một dạng điều kiện xuất hiện trong slide là với mọi \(a\le b\le c\le d\),
\[
  w(a,d)+w(b,c)\ge w(a,c)+w(b,d).
\]
Đây là bất đẳng thức Monge theo quy ước cực tiểu. Nó tương đương với việc
hiệu rời rạc
\[
  w(x,i+1)-w(x,i)
\]
không tăng hoặc không giảm theo \(x\), tùy cách đặt chỉ số và chiều tối ưu.
Trực giác là các đường chi phí của các quyết định cắt nhau có trật tự: quyết
định có chỉ số lớn hơn, nếu đã tốt hơn ở một vị trí, sẽ tiếp tục tốt hơn ở
các vị trí xa hơn.

Từ điều kiện trên suy ra tính đơn điệu của quyết định tối ưu:
\[
  x<y \quad\Longrightarrow\quad \opt(x)\le \opt(y).
\]
Tuy nhiên bản trình chiếu không dừng ở chia để trị theo quyết định đơn điệu;
nó còn xét cách lưu toàn bộ các khoảng mà mỗi ứng viên là tốt nhất.

\section{Biên của ứng viên}

Với một ứng viên \(i\), đặt \(B[i]\) là vị trí nhỏ nhất mà từ đó \(i\) thắng
mọi ứng viên đứng trước nó:
\[
  B[i]=\min\{x: t(i,x)<t(j,x)\ \text{với ứng viên }j<i\ \text{đang xét}\}.
\]
Nếu không tồn tại vị trí như vậy, \(i\) bị loại. Nếu tồn tại, \(i\) chỉ có
thể có ích từ \(B[i]\) trở đi.

Giả sử danh sách ứng viên hiện tại là
\[
  i_1<i_2<\cdots<i_k.
\]
Dạng tăng cường của tính đơn điệu lồi hoàn toàn bảo đảm các biên cũng có thứ
tự:
\[
  B[i_1]<B[i_2]<\cdots<B[i_k].
\]
Do đó khi tính \(f[x]\), ta chỉ cần lấy ứng viên cuối cùng trong danh sách có
biên không vượt quá \(x\):
\[
  j=\max\{r: B[i_r]\le x\},\qquad
  f[x]=t(i_j,x).
\]

\section{Thêm một ứng viên mới}

Khi đã tính xong \(f[x]\), vị trí \(x\) có thể trở thành ứng viên cho các
trạng thái sau. Ta so sánh nó với ứng viên cuối \(i_k\) trong danh sách bằng
cách tìm
\[
  y=\min\{y: t(x,y)<t(i_k,y)\}.
\]
Nếu không tồn tại \(y\), ứng viên \(x\) vô dụng. Nếu \(y\le B[i_k]\), ứng
viên \(x\) thắng \(i_k\) ngay từ trước vùng mà \(i_k\) có ích; do đó \(i_k\)
bị loại, và ta tiếp tục so sánh \(x\) với ứng viên cuối mới. Nếu \(y>B[i_k]\),
ta đặt \(B[x]=y\) và thêm \(x\) vào cuối danh sách.

Quá trình này giống duy trì bao dưới của một họ hàm. Mỗi ứng viên được thêm
một lần và bị xóa nhiều nhất một lần, nên số thao tác danh sách là tuyến
tính. Tổng thời gian phụ thuộc vào chi phí tìm điểm giao nhỏ nhất \(y\):
\begin{itemize}
  \item nếu tính được \(y\) trong \(\Oh(1)\), toàn bộ DP chạy \(\Oh(n)\);
  \item nếu phải tìm nhị phân trên miền \(x\), thời gian là \(\Oh(n\log n)\);
  \item nếu mỗi lần tìm giao mất \(T\), thời gian tổng quát là \(\Oh(nT)\).
\end{itemize}

\section{Dạng bài Two Sawmills}

Một ví dụ được nhắc trong slide là CEOI 2004 \textit{Two Sawmills}. Bài có
\(N\) làng hoặc điểm sản xuất trên một đường, mỗi điểm có trọng lượng gỗ và
vị trí. Sau khi biến đổi bằng tổng tiền tố, một lớp trạng thái có dạng
\[
  d[k,x]=\min_{i<x}\{d[k-1,i]+\text{chi phí chuyển gỗ từ đoạn }(i,x]\}.
\]
Chi phí đoạn có thể viết bằng các tổng tiền tố \(S_x\), \(T_x\), \(P_x\), nên
so sánh hai quyết định \(i\) và \(j\) dẫn tới một ngưỡng \(g(i,j)\). Các công
thức trích xuất được từ slide có các đại lượng \(U_i=d[k-1,i]\), \(h(x)=S_x\)
và các ngưỡng
\[
  g(i_0,i_1)<g(i_1,i_2)<\cdots<g(i_{l-1},i_l).
\]
Đây chính là dạng biên \(B[\cdot]\) viết theo biến đã biến đổi. Khi
\(h(x)\) tăng đơn điệu theo \(x\), ta di chuyển con trỏ trong danh sách ứng
viên và lấy quyết định tương ứng trong \(\Oh(1)\) khấu hao.

Với cách tìm giao trực tiếp, slide ghi nhận độ phức tạp có thể giảm từ
\(\Oh(kn\log n)\) xuống \(\Oh(kn)\) cho \(k\) lớp DP, tùy công thức chi phí
cụ thể.

\section{Quan hệ với bao lồi}

Ở một số bài, hiệu giữa hai ứng viên có dạng tuyến tính:
\[
  f_2(x)-f_1(x)=kx+b.
\]
Khi đó mỗi ứng viên tương ứng với một đường thẳng \(y=k_ix+b_i\), còn việc
tìm giá trị nhỏ nhất là truy vấn bao dưới. Nếu các hệ số góc hoặc điểm truy
vấn có thứ tự đơn điệu, ta dùng hàng đợi lồi để đạt \(\Oh(n)\). Nếu không có
đủ thứ tự, ta dùng cây cân bằng hoặc Li Chao tree để đạt
\(\Oh((n+q)\log n)\).

Điểm khác của bài trình chiếu là nó không chỉ xử lý đường thẳng. Các hàm
\(t(i,x)\) có thể là hàm rời rạc phức tạp hơn, miễn là tính chất Monge mạnh
đủ để bảo đảm thứ tự điểm giao và cho phép loại ứng viên bằng biên \(B[i]\).

\section{Cách dùng trong thi}

Khi gặp một DP dạng \(\min_i f[i]+w(i,x)\), nên kiểm tra ba câu hỏi:
\begin{enumerate}
  \item \(w\) có thỏa bất đẳng thức tứ giác/Monge không?
  \item quyết định tối ưu có đơn điệu theo trạng thái không?
  \item điểm mà một quyết định mới thắng quyết định cũ có tìm được nhanh
  không?
\end{enumerate}
Nếu chỉ có câu 1 và 2, chia để trị tối ưu DP thường đủ để giảm về
\(\Oh(n\log n)\) cho một lớp. Nếu câu 3 cũng có câu trả lời tốt, danh sách
ứng viên với biên \(B[i]\) có thể giảm tiếp về tuyến tính.

\section{Kết luận}

Thông điệp của slide là tính đơn điệu lồi hoàn toàn không chỉ nói rằng vị trí
tối ưu tăng dần. Trong nhiều bài, ta còn quản lý được toàn bộ thứ tự các
ứng viên và các biên nơi chúng bắt đầu tốt hơn nhau. Khi nhận ra cấu trúc
này, một DP bậc hai có thể được biến thành một thuật toán tuyến tính hoặc
gần tuyến tính mà vẫn giữ chứng minh rõ ràng dựa trên bất đẳng thức Monge.

\end{document}
